\begin{problem}[第35届全国中学生物理竞赛决赛][负折射介质光学]{101}
    介质的折射率 $n$ 可以大于 0，也可以小于 0。$n$ 小于 0 的介质称为负折射介质。光在负折射介质内传播，其光程为负值（相位随传播距离的变化规律与在折射率为正的介质中的相反）。如果定义折射角与入射角在界面法线同侧时折射角为负，可以证明折射定律在介面两边有负折射介质时仍然成立，即 $n_{1}\sin\theta_{1}=n_{2}\sin\theta_{2}$ ，其中的 $n_{1}$ 和 $n_{2}$ 均可以大于 0 或小于 0， $\theta_{2}$ 为折射角。
    \begin{subproblem}
        设想一束平行光入射到界面上，根据惠更斯原理，在答题纸上画出图a和图b所示情况下进入介质2的光线及对应的子波的示意图，并依此证明折射定律成立；
    \end{subproblem}
    \begin{subproblem}
        如图 c 所示，半径为 $R$ 的球面将空间隔开为两个区域，其折射率分别记为 $n_{1}$ （ $n_{1}>0$ ）、 $n_{2}$ （ $n_{2}<0$ ），C 点是球面的球心，取某一光轴与球面的交点 O 为原点。图中已画出此情形下一段入射光线和折射光线，$x$ 和 $y$ 分别为入射光线、折射光线与光轴的交点坐标。记物距为 $s_{1}$ ，像距为 $s_{2}$ 。在傍轴近似下导出球面的成像公式和横向放大率公式。请明确指出最后结果中各个量的正负号约定。
    \end{subproblem}
    \begin{subproblem}
        设介质 1 为空气，即 $n_{1} \approx 1$ ， $n_{2}$ 可大于 0 也可小于 0。在球面（参考图 c）前放置一普通薄凸透镜，透镜的光轴通过球心 C，焦点位于负折射介质区域内，透镜的焦距 $f = 1.5R$，透镜中心 $O'$ 点与 O 点的距离为 $d$。一束沿光轴传播的平行光入射到薄透镜。分别就表中四组参数计算入射光在光轴上会聚点离 O 点的距离，并在答题纸上画出序号 4 情形的光路示意图。

        \begin{tabular}{|c|c|c|}
            \hline
            序号 & $n_2$ & $d$ \\
            \hline
            1 & 1.5 & 0.35R \\
            \hline
            2 & 1.5 & 0.85R \\
            \hline
            3 & -1.5 & 0.35R \\
            \hline
            4 & -1.5 & 0.85R \\
            \hline
        \end{tabular}
    \end{subproblem}
\end{problem}